\scanpage{111}
% archive foliation: 110
comme produit $U^{\mathrm{an}}\times P_{x_0}^{\mathrm{an}}$, donc $\sigma|U$ s'identifie à une application $U^{\mathrm{an}}\longrightarrow P_{x_0}^{\mathrm{an}}$, dont les composées avec les iso. $P_{x_0}^{\mathrm{an}}\longrightarrow\mathbb P^1(\mathbb C)$ donnent les germes d'applications \deleted{\unsure{analytiques}} du faisceau $\mathcal U$.

Quelle est l'interprétation des sections horizontales \inserted{locales} de $P^{\mathrm{an}}$ dans la description \textcircled{e} ? \unsure{Pour cela on a un} fibré principal (torseur) $\mathcal T$ sur $X$ [\unsure{donc} sur $X^{\mathrm{an}}$ ; il faut prendre les sections horizontales de $\mathcal T^{\mathrm{an}}=\underline{\operatorname{Isom}}(\mathbb P^1(\mathbb C)_{X^{\mathrm{an}}},P)$ et les appliquer aux pts de $\mathbb P^1(\mathbb C)$ pour trouver sections horizontales de $P$ ; la section horizontale qui passe par $\unsure{y_0}\in P_{x_0}$ s'obtient en choisissant un \deleted{\unsure{point}} élément $f_{x_0}$ de $\mathcal T_{x_0}^{\mathrm{an}}$, d'où \deleted{\unsure{le}} le pt $p_0$ de $\mathbb P^1(\mathbb C)_{X^{\mathrm{an}}}$ [\unsure{car} $f_{x_0}$ définit un iso $\mathbb P^1(\mathbb C)_{\deleted{\unsure{x_0}}}\xrightarrow{\sim}P_x(\ni p)$], et prendre $x\mapsto\widetilde f_{x_0}(x)$, où $\widetilde f_{x_0}$ est la section horizontale \inserted{locale} de $P^{\mathrm{an}}$ passant par $f_{x_0}$. Or]

$P$ s'interprète comme $\mathcal T/B$ ($B\subset\mathrm{GP}(1)_{\mathbb C}$, Borel). Les sections \deleted{\unsure{les}} horizontales \inserted{de $P^{\mathrm{an}}$} sont les sections horiz. de $\mathcal T^{\mathrm{an}}$ mod action de $B^{\mathrm{an}}\deleted{\unsure{\mathbb C}}$ -- i.e. sections de $\mathcal U_{\deleted{\unsure{\mathcal T^{\mathrm{an}}}}}\subset\underline{\operatorname{Appl}}(X^{\mathrm{an}},\mathbb P^{1\mathrm{an}})$, mod action de $B$, i.e. sections de $\mathcal U/B^{\mathrm{an}}$. [Interpréter une section de $P$ disjointe de $\sigma$ comme une connexion \inserted{$D$} sur $\underline\omega$, \inserted{$D:\underline\omega\longrightarrow\underline\omega^2$} on doit prendre les connexions $D$ telles que $\underline c_D=c_0$ (la connexion projective donnée).

Quelles sont les sections \inserted{(non horizontales)} de $\mathcal T/B_{\unsure{u}}$, et quel est
